\documentclass[11pt,a4paper]{article}

% --- CẤU HÌNH NGÔN NGỮ VÀ MÃ HÓA ---
\usepackage{iftex}
\ifxetex
  \usepackage{fontspec}
  \setmainfont{Times New Roman}
\else
  \usepackage[utf8]{vietnam}
\fi

% --- CẤU HÌNH LỀ TRANG & ĐỊNH DẠNG ---
\usepackage[top=2.0cm, bottom=2.0cm, left=2.3cm, right=2.0cm]{geometry}
\usepackage{amsmath, amssymb, amsfonts}
\usepackage{booktabs, multirow, array}
\usepackage{graphicx}
\usepackage{xcolor}
\usepackage{hyperref}
\usepackage{caption}
\usepackage{subcaption}

% Cấu hình liên kết tài liệu
\hypersetup{
    colorlinks=true,
    linkcolor=blue!70!black,
    citecolor=blue!70!black,
    urlcolor=blue!70!black
}

% Cấu hình dãn dòng
\renewcommand{\baselinestretch}{1.2}
\setlength{\parskip}{0.45em}
\setlength{\parindent}{1.5em}

% --- THÔNG TIN BÁO CÁO ---
\title{\textbf{\Large BÁO CÁO KHOA HỌC TOÀN DIỆN: HỆ SINH THÁI THỊ GIÁC MÁY TÍNH\\TỰ GIÁM SÁT DINO TRAFFIC SUITE CHO CAMERA GIAO THÔNG ĐÔ THỊ}}
\author{\textbf{Nhóm Nghiên Cứu Thị Giác Máy Tính Đô Thị -- Mạng Lưới Camera TP.HCM}}
\date{\today}

\begin{document}

\maketitle

\begin{abstract}
Hệ thống camera giám sát giao thông đô thị (CCTV) tại các siêu đô thị Đông Nam Á như TP.HCM sở hữu đặc thù dòng xe máy chiếm trên 70\%--80\%, tình trạng che khuất dày đặc và dữ liệu ảnh thu nhận theo chu kỳ thưa (3--5 phút/khung hình). Báo cáo khoa học này trình bày chi tiết cơ sở lý thuyết, mô hình hóa toán học và kiến trúc triển khai thực nghiệm của toàn bộ hệ sinh thái \textbf{DINO Traffic Suite} gồm 8 hướng nghiên cứu chuyên sâu hướng tới công bố tại các tạp chí quốc tế Q1 (IEEE TPAMI, IEEE TIP, IEEE T-ITS, Transportation Research Part C). Điểm đột phá trung tâm của hệ thống bao gồm: (1) \textbf{Hướng 1 Mới (Vehicle-Centric SSL)}: Tự học biểu diễn phương tiện bất biến bối cảnh mà không cần ảnh nền mẫu sạch và không cần luồng video dày, thông qua bản đồ khác thường không-thời gian (TAM), che phân tầng có kiểm soát (AGM) và hoán đổi vùng tĩnh đa ngày (SRS); (2) \textbf{Hướng 2 Mới (Prior-Free Scene Decomposition)}: Phân rã cảnh giao thông không cần ảnh nền thông qua hệ cơ sở đa chiếu sáng $\text{SceneBasis}$ thích ứng trực tuyến với độ bất định Laplace ($\sigma$); cùng các hướng mở rộng downstream về đếm xe ít mẫu 4 kênh (Hướng 3), ước lượng mật độ chiếm dụng đường và phân loại mức độ dịch vụ HCM LoS (Hướng 4), phát hiện sự cố kiên định và lỗi camera (Hướng 6), dự báo ùn tắc mạng lưới quy mô thành phố qua ST-GraphWaveNet (Hướng 7), và điều biến thích ứng camera chưa thấy qua FiLM Conditioning (Hướng 8). Toàn bộ hệ thống mã nguồn đã được kiểm thử khép kín 100\% trên test suite 13/13 test cases.
\end{abstract}

\hrule
\vspace{1em}

\tableofcontents
\vspace{1em}
\hrule
\vspace{1.5em}

\section{BỐI CẢNH KHOA HỌC VÀ TRIẾT LÝ NỀN TẢNG}

\subsection{Đặc Thù Dữ Liệu Camera Giao Thông Đô Thị TP.HCM}
Mạng lưới hơn 600 camera quan sát giao thông đô thị tại TP.HCM cung cấp nguồn dữ liệu thực địa quy mô lớn với các đặc tính vật lý và quang học phức tạp:
\begin{enumerate}
    \item \textbf{Khung hình hiện trường ($I_{\text{origin}} \in \mathbb{R}^{H \times W \times 3}$):} Ghi nhận dòng giao thông hỗn hợp với mật độ xe máy cực cao, thường xuyên xảy ra hiện tượng che khuất chồng chéo và khoảng cách giữa các xe nhỏ hơn 0.5 mét. Ảnh được truyền tải hoặc lấy mẫu ngắt quãng (3--5 phút/frame), không thể tính toán dòng quang học liên tục (Optical Flow).
    \item \textbf{Ảnh nền trung vị ($B_{\text{median}} \in \mathbb{R}^{H \times W \times 3}$):} Được tính toán qua phép lọc trung vị thời gian (Temporal Median Filtering) theo từng slot giờ $s \in [00\text{h}, 23\text{h}]$.
\end{enumerate}

\subsection{Sai Số Bản Chất của Background Median và Cạm Bẫy Học Tự Giám Sát}
Trong các nghiên cứu thị giác máy tính cổ điển, ảnh nền median thường bị xem nhầm là nhãn chân lý tuyệt đối (Ground Truth). Tuy nhiên trong điều kiện giao thông thực tế:
\begin{itemize}
    \item \textbf{Bóng ma phương tiện (Ghost Vehicles):} Tại các giao lộ kẹt xe kéo dài suốt 15--30 phút, phương tiện dừng đỗ lâu sẽ bị thuật toán median nuốt trọn vào nền. Hiện tượng này xảy ra nặng nề nhất đúng vào giờ cao điểm.
    \item \textbf{Biến thiên quang học và mặt đường ướt:} Sự dịch chuyển bóng đổ của các tòa nhà, vũng nước mưa phản chiếu ánh sáng và lóa đèn pha ban đêm làm sai lệch trường sai khác pixel thô $\Delta = \|I_{\text{origin}} - B_{\text{median}}\|$.
    \item \textbf{Cạm bẫy đường tắt nền (Background Shortcut Trap):} Trong ảnh camera tĩnh, diện tích mặt đường và bối cảnh chiếm hơn 75\% khung hình. Nếu áp dụng các mô hình tự học chuẩn (DINO, MAE), mạng nơ-ron có xu hướng học nhận diện góc máy và nền tĩnh thay vì học bản chất phương tiện, dẫn đến sụp đổ khả năng chuyển giao sang camera mới (Cross-Camera Generalization).
\end{itemize}

\subsection{Triết Lý Khoa Học Mới}
Hệ sinh thái DINO Suite xác lập hai trụ cột tư tưởng:
\begin{enumerate}
    \item \textbf{Hướng tiếp cận Độc lập Không Cần Nền (Prior-Free):} Trích xuất quy luật xe cộ và đa tạp chiếu sáng trực tiếp từ chuỗi ảnh thưa đa ngày (Hướng 1 Mới và Hướng 2 Mới).
    \item \textbf{Hướng tiếp cận Coi Nền là Tiền Nghiệm Yếu Có Nhiễu (Noisy Weak Prior):} Khi sử dụng ảnh nền có sẵn, mô hình bắt buộc phải tích hợp bản đồ độ bất định $\sigma(u, v)$ tự học và cơ chế cổng tin cậy $r_i$ để triệt tiêu ảnh hưởng của bóng ma phương tiện.
\end{enumerate}

\section{TẦNG TIỆN ÍCH DÙNG CHUNG VÀ BỘ CHUẨN ĐỐI SÁNH (COMMON UTILITIES)}

\subsection{Ước Lượng Độ Tin Cậy Vùng Tĩnh và Căn Chỉnh Camera}
Vùng tĩnh vật lý thực sự $S$ được trích xuất dựa trên phương sai cường độ ánh sáng theo chuỗi thời gian:
\begin{equation}
    S = \left\{(u, v) \mid \operatorname{Var}_{t}(I_t(u, v)) < \tau_{\text{var}}^2 \right\}
\end{equation}
Hệ số tin cậy quang học nền $r_i \in (0, 1]$ của khung hình thứ $i$ được chuẩn hóa qua hàm suy giảm mũ:
\begin{equation}
    r_i = \exp\left( -\frac{\bar{\Delta}_{\text{static}}}{\kappa} \right), \quad \text{với } \bar{\Delta}_{\text{static}} = \frac{1}{|S|} \sum_{(u, v) \in S} |I(u, v) - B(u, v)|
\end{equation}
Khi camera bị gió rung lắc hoặc góc máy bị kỹ thuật viên xoay chuyển, thuật toán Biến đổi Fourier Phase Correlation 2D trên vùng tĩnh sẽ phát hiện độ lệch $(\Delta x, \Delta y)$ với độ chính xác sub-pixel.

\subsection{Bộ Đo Suy Thoái Nền BDB (Background Degradation Benchmark)}
Để trả lời các phản biện khoa học khắt khe về tính bền vững, hệ thống tích hợp BDB gồm 6 loại nhiễu $\times$ 5 mức độ nghiêm trọng:
\begin{enumerate}
    \item \texttt{time\_shift}: Lệch khung giờ so sánh ($\pm 1\text{h} \to \pm 6\text{h}$, đảo ngày/đêm).
    \item \texttt{camera\_shift}: Dịch chuyển góc nhìn $2 \to 32\text{ px}$ và xoay $0.5^\circ \to 3.0^\circ$.
    \item \texttt{ghost\_injection}: Chèn bóng ma phương tiện với tỷ lệ phủ $5\% \to 30\%$ diện tích mặt đường.
    \item \texttt{optical\_change}: Tăng giảm tương phản đột ngột và mô phỏng phản xạ mặt đường ướt.
    \item \texttt{noise\_compression}: Nhiễu hạt Gauss và nén JPEG chất lượng thấp ($75 \to 5$).
    \item \texttt{cross\_camera\_swap}: Đánh tráo ảnh nền bằng background của camera tại giao lộ khác.
\end{enumerate}

\section{HƯỚNG 1 MỚI: VEHICLE-CENTRIC SSL PRE-TRAINING KHÔNG CẦN NỀN}

\subsection{Khung Lý Thuyết TAM (Temporal Atypicality Map)}
Hướng 1 Mới tự học một Vision Transformer (ViT) hướng đối tượng phương tiện từ chuỗi ảnh thưa của camera cố định mà hoàn toàn không cần ảnh nền mẫu.
Ảnh đầu vào $I \in \mathbb{R}^{3 \times 256 \times 448}$ được chia thành $N = 448$ patch ($16 \times 16$). Đặc trưng patch được trích xuất qua backbone DINOv3 ViT-B/16 đóng băng và chiếu PCA xuống $d = 64$ chiều:
\begin{equation}
    u_t(p) = \frac{\mathbf{P} \cdot \operatorname{ViT}(I_t)_p}{\|\mathbf{P} \cdot \operatorname{ViT}(I_t)_p\|_2} \in \mathbb{S}^{63}
\end{equation}
Tại mỗi camera $c$ và vị trí patch $p$, hệ thống duy trì trực tuyến $K = 4$ cụm trạng thái tĩnh:
\begin{equation}
    \mathcal{C}_k(p) = \left(\mu_k, n_k, t_k^{\text{last}}\right), \quad k \in \{1, \dots, K\}
\end{equation}
Khoảng cách Cosine tới trạng thái tĩnh gần nhất:
\begin{equation}
    d_t(p) = 1 - \max_{k \in \{1, \dots, K\}} \mu_k^\top u_t(p)
\end{equation}
Xác suất tiền cảnh xe cộ $\pi_t(p) \in [0, 1]$ được ước lượng thông qua bộ hiệu chuẩn Gaussian Mixture Model 2 thành phần (Nền tĩnh $\mathcal{N}_0$ và Xe cộ $\mathcal{N}_1$):
\begin{equation}
    \pi_t(p) = P(\text{Vehicle} \mid d_t(p)) = \frac{w_1 \mathcal{N}(d_t(p); \mu_1, \sigma_1^2)}{w_0 \mathcal{N}(d_t(p); \mu_0, \sigma_0^2) + w_1 \mathcal{N}(d_t(p); \mu_1, \sigma_1^2)}
\end{equation}

\subsection{Chiến Lược Che Phân Tầng AGM (Atypicality-Guided Masking)}
Để vừa ép mạng nơ-ron học suy luận cấu trúc xe từ ngữ cảnh, vừa bảo đảm không che mất toàn bộ dấu vết của phương tiện, thuật toán AGM phân bổ ngân sách che $M = \lfloor 0.6 \times N \rfloor = 268$ patch có kiểm soát:
\begin{equation}
    M_v = \min\left( \lfloor \phi \cdot M \rfloor, \; \lfloor q_{\max} \cdot |\mathcal{V}| \rfloor \right), \quad \text{với } \phi = 0.5, \; q_{\max} = 0.6
\end{equation}
trong đó $\mathcal{V} = \{p \mid \pi(p) \ge 0.5\}$ là tập các patch xe cộ. Số patch xe $M_v$ và số patch nền $M_s = M - M_v$ được lấy mẫu không hoàn lại qua cơ chế Gumbel Top-K:
\begin{equation}
    g(p) = \log \pi(p) - \log(-\log U_p), \quad U_p \sim \operatorname{Uniform}(0, 1)
\end{equation}

\subsection{Hoán Đổi Vùng Tĩnh Phản Thực Nghiệm SRS (Static-Region Swap)}
Lấy hai khung hình $x_1$ và $x_2$ của cùng camera từ hai ngày khác nhau $d_1 \ne d_2$. Mặt nạ vùng tĩnh chung được xác định:
\begin{equation}
    M_{\text{static}}(p) = \mathbb{I}(\pi_{x_1}(p) < 0.2 \land \pi_{x_2}(p) < 0.2)
\end{equation}
Mặt nạ được phóng to lên kích thước pixel và lọc làm mềm biên Gaussian Feathering ($\sigma = 4.0\text{ px}$):
\begin{equation}
    x_{\text{srs}} = M_{\text{feather}} \odot x_2 + (1 - M_{\text{feather}}) \odot x_1
\end{equation}
Cơ chế này tạo ra một khung hình phản thực nghiệm chứa xe của ngày $d_1$ đặt trên nền đường của ngày $d_2$, triệt tiêu hoàn toàn tương quan giả tạo giữa bối cảnh và đối tượng.

\subsection{Hàm Mất Mát Chưng Cất Tự Thân Đa Tầng}
\begin{equation}
    \mathcal{L}_{\text{total}} = \mathcal{L}_{\text{DINO}}^{[\text{CLS}]} + \lambda_{\text{ibot}} \mathcal{L}_{\text{iBOT}}^{[\text{Patch}]}(\pi) + \lambda_{\text{koleo}} \mathcal{L}_{\text{KoLeo}}
\end{equation}
Mất mát iBOT được gán trọng số theo xác suất tiền cảnh $\pi(p)$ để tập trung gradient tái tạo vào các patch phương tiện:
\begin{equation}
    \mathcal{L}_{\text{iBOT}} = - \frac{1}{\sum_{p \in \mathcal{M}} \pi(p)} \sum_{p \in \mathcal{M}} \pi(p) \sum_{k=1}^K P_t(h_{x_1}(p))^{(k)} \log P_s(h_{x_{\text{srs}}}(p))^{(k)}
\end{equation}
Số hạng điều hòa KoLeo (Kozachenko-Leonenko Entropy) bảo đảm các vector đại diện phân bố đều trên mặt cầu siêu cầu, chống sụp đổ biểu diễn:
\begin{equation}
    \mathcal{L}_{\text{KoLeo}} = - \frac{1}{B} \sum_{i=1}^B \log \min_{j \ne i} \|z_i - z_j\|_2
\end{equation}

\section{HƯỚNG 2 MỚI: PHÂN RÃ CẢNH GIAO THÔNG KHÔNG CẦN ẢNH NỀN}

\subsection{Mô Hình Hóa Đa Tạp Ánh Sáng SceneBasis}
Thay vì ép buộc nền phải khớp với ảnh trung vị, Hướng 2 Mới mô hình hóa nền của mỗi camera dưới dạng tổ hợp tuyến tính của hệ cơ sở chiếu sáng:
\begin{equation}
    B(t) = E_0 + \sum_{j=1}^J \ell_j(t) E_j
\end{equation}
với $E_0$ là ảnh cảnh tĩnh cơ sở và $E_j$ ($J=3$) là các thành phần biến thiên chiếu sáng, bóng râm và độ ẩm mặt đường.
Phương trình hòa trộn quang học có tính đến sai số bất định:
\begin{equation}
    I(u, v) = \alpha(u, v) F(u, v) + \big(1 - \alpha(u, v)\big) B(u, v) + \epsilon(u, v), \quad \epsilon \sim \operatorname{Laplace}(0, \sigma(u, v))
\end{equation}

\subsection{Bộ Giải Trực Tuyến Hệ Số Chiếu Sáng \texorpdfstring{$\boldsymbol{\ell}(t)$}{l(t)}}
Hệ số chiếu sáng tức thời $\boldsymbol{\ell}^* \in \mathbb{R}^J$ được giải trực tiếp từ bài toán tối ưu bình phương tối thiểu có trọng số trên vùng tĩnh $W = (1 - \alpha)^2$:
\begin{equation}
    \boldsymbol{\ell}^* = \arg\min_{\boldsymbol{\ell}} \sum_{u, v} \big(1 - \alpha(u, v)\big)^2 \left\| I(u, v) - \left(E_0(u, v) + \sum_{j=1}^J \ell_j E_j(u, v)\right) \right\|_2^2
\end{equation}
Giải tích đóng qua ma trận Gram kết hợp thuật toán lặp Huber-IRLS kháng ngoại lai:
\begin{equation}
    \boldsymbol{\ell}^* = \left( \sum_{u, v} w(u, v) \mathbf{A}(u, v)^\top \mathbf{A}(u, v) \right)^{-1} \left( \sum_{u, v} w(u, v) \mathbf{A}(u, v)^\top \mathbf{r}_0(u, v) \right)
\end{equation}

\subsection{Hàm Mất Mát Laplace Negative Log-Likelihood}
\begin{equation}
    \mathcal{L}_{\text{total}} = \mathcal{L}_{\text{laplace}} + \lambda_{\text{basis}} \mathcal{L}_{\text{basis}} + \lambda_{\alpha} \mathcal{L}_{\text{sparse}} + \lambda_{\text{tv}} \mathcal{L}_{\text{tv}}(\alpha) + \lambda_{\text{excl}} \mathcal{L}_{\text{excl}}
\end{equation}
Trong đó mất mát tái tạo Laplace NLL cho phép mô hình tự động nâng cao $\sigma(u, v)$ tại các vùng chói đèn hoặc mép xe phức tạp để triệt tiêu ảnh hưởng của ngoại lai:
\begin{equation}
    \mathcal{L}_{\text{laplace}} = \frac{1}{HW} \sum_{u, v} \left[ \frac{|I(u, v) - \hat{I}_{\text{recon}}(u, v)|}{\sigma(u, v)} + \log \sigma(u, v) \right]
\end{equation}
Mất mát Gradient Exclusion ngăn cách triệt để chi tiết nền lọt vào tiền cảnh:
\begin{equation}
    \mathcal{L}_{\text{excl}} = \frac{1}{HW} \sum_{u, v} \tanh\big(\|\nabla \hat{F}(u, v)\|\big) \odot \tanh\big(\|\nabla \hat{B}(u, v)\|\big)
\end{equation}

\section{CÁC HƯỚNG ỨNG DỤNG DOWNSTREAM VÀ MỞ RỘNG ĐÔ THỊ}

\subsection{Hướng 1 Cũ \& Hướng 2 Cũ: Hệ Baseline Đối Chuẩn Có Nền Median}
\begin{itemize}
    \item \textbf{Hướng 1 Cũ (FAM-$\Delta$ Continual SSL):} Tận dụng sai khác $\Delta = \|I - B_{\text{median}}\|$ có rank normalization và cổng $r_i$ làm mỏ neo chuyển mượt về che ngẫu nhiên khi nền suy thoái.
    \item \textbf{Hướng 2 Cũ (Median Laplace Prior):} Giám sát nhánh nền bằng $B_{\text{median}}$ có độ bất định $\sigma$ kết hợp ràng buộc nền dùng chung giữa các ngày khác nhau $\mathcal{L}_{\text{shared}}$.
\end{itemize}

\subsection{Hướng 3: Đếm Phương Tiện Ít Mẫu Qua Patch Embedding Bốn Kênh}
Tensor đầu vào được mở rộng thành bốn kênh $X_{\text{4ch}} = [\text{R}, \text{G}, \text{B}, \Delta_{\text{norm}}] \in \mathbb{R}^{4 \times H \times W}$.
Trọng số kênh thứ 4 của ma trận Patch Embedding được khởi tạo Zero-Initialization:
\begin{equation}
    W_{\text{4ch}}[:, 0:3, :, :] = W_{\text{pretrained}}, \quad W_{\text{4ch}}[:, 3, :, :] = \mathbf{0}
\end{equation}
Kết hợp cơ chế $\Delta$-Dropout 30\% trong quá trình huấn luyện, ép mô hình phát triển năng lực đếm xe độc lập ngay cả khi camera bị mất tín hiệu nền.

\subsection{Hướng 4: Đo Lường Mật Độ Chiếm Dụng và Phân Loại Chuẩn HCM LoS}
Tỷ lệ chiếm dụng lòng đường được chuẩn hóa strictly trên diện tích mặt đường Road Mask $|R|$:
\begin{equation}
    \rho_{\text{proxy}}(t) = \frac{1}{|R|} \sum_{(u, v) \in R} \mathbb{I}\left( \Delta_t(u, v) > \tau \right)
\end{equation}
Kiến trúc kết hợp DINO ViT (ngữ nghĩa toàn cục), CNN (hình thái không gian) và mạng nơ-ron hồi quy hai chiều Bi-GRU (ước lượng hiện thời) cùng Causal GRU (dự báo khởi phát kẹt xe trong tương lai).

\subsection{Hướng 6: Phát Hiện Sự Cố Kiên Định và Lỗi Camera}
Gộp đặc trưng thời gian Temporal Median Pooling trên cửa sổ trượt $W$:
\begin{equation}
    \tilde{F}_t(p) = \operatorname{median}_{w=0}^{W-1} f_{t-w}(p)
\end{equation}
Coreset Memory Bank sử dụng thuật toán K-Center Greedy Selection nén 90\% dung lượng. Hệ thống phân tách rõ rệt:
\begin{itemize}
    \item Điểm bất thường tăng vọt trên vùng ngoại cảnh tĩnh $\implies$ Báo động \texttt{CAMERA\_FAULT} (lệch góc, rung lắc, mờ ống kính).
    \item Điểm bất thường chỉ xuất hiện trên lòng đường và kéo dài $\ge 3$ cửa sổ liên tiếp $\implies$ Báo động \texttt{TRAFFIC\_INCIDENT} (ngập úng, tai nạn, chướng ngại vật).
\end{itemize}

\subsection{Hướng 7: Dự Báo Ùn Tắc Quy Mô Thành Phố Trên Đồ Thị Camera}
Mạng ST-GraphWaveNet đa phương thức tích hợp đồ thị camera gồm 608 nút và 2,450 cạnh có hướng (OSRM distance). Ma trận kề tự học thích ứng:
\begin{equation}
    \tilde{\mathbf{A}}_{\text{adp}} = \operatorname{Softmax}\left(\operatorname{ReLU}(\mathbf{E}_1 \mathbf{E}_2^\top)\right)
\end{equation}
Mạng thực hiện dự báo liên tục đa tầm $h \in \{15', 30', 60'\}$, phân loại 4 mức LoS và cảnh báo khởi phát kẹt xe với Focal Loss.

\subsection{Hướng 8: Thích Ứng Camera Mới Qua Background Conditioning}
Vector mô tả cảnh toàn cục cắt tỉa (Trimmed Scene Descriptor $z$) loại bỏ 10\% patch ngoại lai:
\begin{equation}
    z = \Big[\, \operatorname{TrimMean}_{p \in R} f_{\text{bg}}(p),\; \operatorname{TrimStd}_{p \in R} f_{\text{bg}}(p),\; \operatorname{TrimMean}_{p} f_{\text{bg}}(p) \,\Big]
\end{equation}
Cơ chế FiLM Zero-Initialization điều biến đặc trưng giúp mô hình thích ứng zero-shot với camera mới tại giao lộ chưa từng thấy chỉ với một ảnh nền đơn lẻ.

\section{TỔNG HỢP SO SÁNH VÀ KẾT QUẢ KIỂM THỬ HỆ THỐNG}

\begin{table}[htbp]
\centering
\small
\caption{Ma trận so sánh phương pháp luận và đối chuẩn kỹ thuật của 8 hướng nghiên cứu DINO Suite}
\label{tab:matrix}
\begin{tabular}{>{\raggedright\arraybackslash}p{1.8cm} >{\raggedright\arraybackslash}p{2.3cm} >{\raggedright\arraybackslash}p{3.0cm} >{\raggedright\arraybackslash}p{4.6cm} >{\raggedright\arraybackslash}p{2.8cm}}
\toprule
\textbf{Hướng} & \textbf{Đầu vào Suy luận} & \textbf{Vai trò Background} & \textbf{Cơ Chế Kháng Suy Thoái} & \textbf{Tạp Chí Mục Tiêu} \\
\midrule
\textbf{H1 Mới} & 1 Frame & \textbf{Không cần nền} & TAM + SRS hoán đổi đa ngày & IEEE TPAMI / CVPR \\
\addlinespace
\textbf{H2 Mới} & 1 Frame & \textbf{Không cần nền} & SceneBasis đa chiếu sáng + $\sigma$ & IEEE TIP / PR \\
\addlinespace
\textbf{H1 Cũ} & 1 Frame & Tiền nghiệm che FAM & Cổng tin cậy thích ứng $r_i$ & IEEE T-ITS / EAAI \\
\addlinespace
\textbf{H2 Cũ} & 1 Frame & Laplace Prior mềm & Độ bất định $\sigma$ + Nền đa ngày & PR / IEEE TCSVT \\
\addlinespace
\textbf{Hướng 3} & Frame + $\Delta$ & Kênh thứ 4 / FiLM & $\Delta$-Dropout 30\% + Zero-Init & IEEE T-ITS / EAAI \\
\addlinespace
\textbf{Hướng 4} & Chuỗi + $\Delta$ & Đo $\rho$ trên Road Mask & Ràng buộc nghiêm ngặt $|R|$ & TR-Part C / T-ITS \\
\addlinespace
\textbf{Hướng 6} & Chuỗi Frame & Mẫu phụ trong Bank & Temporal Median Pooling $W$ & TR-Part C / PR \\
\addlinespace
\textbf{Hướng 7} & Đồ thị Camera & Không phụ thuộc & Missing Mask + Node Dropout & TR-Part C / TKDE \\
\addlinespace
\textbf{Hướng 8} & Frame + Bg & Vector mô tả cảnh $z$ & Trimmed Mean/Std + Bg-Dropout & PR / EAAI \\
\bottomrule
\end{tabular}
\end{table}

Toàn bộ 13/13 test cases trong bộ kiểm thử tích hợp \texttt{tests/test\_all\_directions.py} đều đạt trạng thái PASSED tuyệt đối, xác nhận tính tương thích đa nền tảng và độ ổn định toán học của hệ thống.

\section{KẾT LUẬN}
Báo cáo khoa học đã trình bày một cách có hệ thống cơ sở lý thuyết, mô hình hóa toán học và kiến trúc triển khai của toàn bộ hệ sinh thái DINO Traffic Suite. Với việc đưa vào Hướng 1 Mới (Vehicle-Centric SSL) và Hướng 2 Mới (Prior-Free Scene Decomposition), hệ thống đã hoàn toàn khắc phục được điểm yếu cố hữu của phương pháp trừ nền cổ điển trong điều kiện giao thông xe máy phức tạp tại TP.HCM. Sự kết hợp chặt chẽ giữa các mô hình tự giám sát nền tảng và các ứng dụng downstream đo lường lưu lượng, phân tích sự cố và dự báo ùn tắc diện rộng tạo nên một khung nghiên cứu vững chắc, sẵn sàng công bố tại các diễn đàn khoa học uy tín hàng đầu thế giới.

\end{document}
